\documentclass[10pt, a4paper]{article}
\usepackage{preamble}

\title{Statistical Inference II \\
    \large Problem Class}
\author{Luke Phillips}
\date{October 2025}

\begin{document}

\maketitle

\newpage

\tableofcontents

\newpage

\section{Problems Class \texorpdfstring{$1$}{}}

\begin{problem}[Sheet $1$ Question $3$]
    Suppose that $Y$ and $Z$ are independent random variables,
    where $Y$ has a chi-square distribution with $n$ df,
    and $Z$ has a standard normal distribution:
    \[
    f_Z(z) = \frac{1}{\sqrt{2\pi}}e ^ {-\frac{1}{2}z ^ 2}\qquad 
    f_Y(y) = \frac{1}{2 ^ {n / 2}\Gamma(n / 2)}y ^ {n / 2 - 1}e ^ {-y / 2}.
    \]
    Let $X = \frac{Z}{\sqrt{\frac{Y}{n}}}$ and $W = Y$.

    \begin{enumerate}[label = (\alph*)]
        \item Find the joint pdf of $W$ and $X$.

        \item From its definition above,
        $X$ should have a $t$ distribution with $n$ degrees of freedom.
        By marginalising the joint distribution in part (a),
        deduce the pdf of the $t_n$ distribution is
        \[
        f_{X}(x) = \frac{\Gamma((n + 1) / 2)}{\sqrt{n\pi}\Gamma(n / 2)}\left(1 + \frac{x ^ 2}{n}\right) ^ {-(n + 1) / 2}
        \]
    \end{enumerate}

    \begin{solution}
        \begin{enumerate}[label = (\alph*)]
            \item
            $2$D change of variables:
            Inverses:
            $W = Y$,
            $y = w$
            \[
            x = \frac{z}{\sqrt{y / n}} = \frac{z}{\sqrt{w / n}}
            \]
            \[
            z = x\sqrt{w / n}
            \]
            for $w > 0$,
            $x \in \R$.

            \[
            J = \begin{bmatrix}
                \pd[y]{x} & \pd[y]{w} \\ \pd[z]{x} & \pd[z]{w}
            \end{bmatrix} = \begin{bmatrix}
                0 & 1 \\ \sqrt{w / n} & \frac{x}{2n\sqrt{w / n}}
            \end{bmatrix}.
            \]
            So $\det{J} = -\sqrt{w / n}$,
            $|\det{J}| = \sqrt{w / n}$.

            Apply result:
            \[
            f_{Y, Z}(y, z) = f_{Y}(y)f_{Z}(z)
            \]
            \begin{align*}
                f_{X, W}(x, w) &= f_{Y, Z}(y(x, w), z(x, w))|\det{J}| \\
                &= f_{Y}(y(x, w))f_{Z}(z(x, w))|\det{J}| \\
                &= \frac{1}{2 ^ {n / 2}\Gamma(n / 2)}w ^ {n / 2 - 1}e ^ {-w / 2}\times \frac{1}{\sqrt{2\pi}}e ^ {-1 / 2(x_n\sqrt{w / n}) ^ 2}\times \sqrt{\frac{w}{n}} \\
                &= \underbrace{\frac{1}{2 ^ {(n + 1) / 2}\Gamma(n / 2)\sqrt{n\pi}}}_{c}w ^ {(n + 1) / 2 - 1}e ^ {-\frac{1}{2}(w + x ^ 2w / n)} \\
                &= cw ^ {\frac{n + 1}{2} - 1}e ^ {-\frac{w}{2}\left(1 + \frac{x ^ 2}{n}\right)} && w > 0, x \in \R.
            \end{align*}

            \item Marginalise:
            \begin{align*}
                f_{X}(x) &= \int_{-\infty}^{\infty}f(x, w)\,dw \\
                &= \int_{0}^{\infty}w ^ {(n + 1) / 2 - 1}e ^ {-w\times\overbrace{\frac{1}{2}(1 + x ^ 2 / n)}^{g(x)}}\,dw \\
                &= c\int_{0}^{\infty}w ^ {\frac{n + 1}{2} - 1}e ^ {-wg(x)}\,dw.
            \end{align*}
            Honesty condition:
            recall the gamma pdf,
            $V \sim \mathrm{Ga}(\alpha, \beta)$
            \[
            f_{V}(v) = \frac{\beta ^ {\alpha}}{\Gamma(\alpha)}v ^ {\alpha - 1}e ^ {-\beta v}.
            \]
            After marginalisation,
            the integrand looks like $W \sim \mathrm{Ga}\left(\overbrace{\frac{n + 1}{2}}^{= \alpha}, \overbrace{g(x)}^{= \beta}\right)$
            \begin{align*}
                f_{X}(x) &= c\int_{0}^{w}w ^ {\alpha - 1}e ^ {-\beta w}\,dw \\
                &= c\frac{\Gamma(\alpha)}{\beta ^ {\alpha}}\int_{0}^{\infty}\underbrace{\frac{\beta ^ {\alpha}}{\Gamma(\alpha)}w ^ {\alpha - 1}e ^ {-\beta w}}_{= 1}\,dw \\
                &= c\frac{\Gamma(\alpha)}{\beta ^ {\alpha}} \\
                &= \frac{1}{2 ^ {\frac{n + 1}{2}}\Gamma(n / 2)\sqrt{n\pi}} \times \frac{\Gamma((n + 1) / 2)}{\left(\frac{1}{2}(1 + x / n)\right) ^ {\frac{n + 1}{2}}} \\
                &= \frac{1}{\sqrt{n\pi}}\frac{\Gamma\left(\frac{n + 1}{2}\right)}{\Gamma\left(\frac{n}{2}\right)}\frac{1}{2 ^ {\frac{n + 1}{2}}}\left(\frac{1}{2}\right) ^ {\frac{n + 1}{2}}\left(1 + \frac{x ^ 2}{n}\right) ^ {-\frac{n + 1}{2}} \\
                &= \frac{1}{\sqrt{n\pi}}\frac{\Gamma\left(\frac{n + 1}{2}\right)}{\Gamma\left(\frac{n}{2}\right)}\left(1 + \frac{x ^ 2}{n}\right) ^ {-\frac{n + 1}{2}},\qquad x \in \R.
            \end{align*}
        \end{enumerate}
    \end{solution}
\end{problem}

\begin{problem}[Sheet $3$ Question $4$]
    In a cohort study on the incidence of ischaemic heard disease
    (IHD)
    $337$ male candidates were enrolled.
    Each man was categorised as non-exposed
    (group $1$)
    or exposed
    (group $2$)
    based on their daily energy consumption to summarise his average level of physical activity.
    
    For each group,
    the number of person years were $Y_1 = 2768.9$ and $Y_2 = 1857.5$,
    and the number of IHD cases $D_1 = 17$ and $D_2 = 28$ was registered thereafter.

    Given $Y_i$,
    we assume that the $D_i$ are IID $\Po(\lambda_i, Y_i$,
    $i = 1, 2$,
    where $\lambda > 0$ is the group-specific incidence rate.

    \begin{enumerate}[label = (\alph*)]
        \item Separately for each group,
        derive the MLE $\hat{\lambda}_i$,
        find the observed information,
        and evaluate the estimates given the data.

        \item In order to analyse whether $\lambda_1 = \lambda_2$,
        we re-parametrise the model with $\lambda = \lambda_1$ and $\theta = \frac{\lambda_2}{\lambda_1}$.
        Show that the joint log-likelihood of $\lambda$ and $\theta$ has the following form:
        \[
        \mathcal{L}(\lambda, \theta) = c + D\log{\lambda} + D_2\log{\theta} - \lambda Y_1 - \theta\lambda Y_2,
        \]
        where $D = D_1 + D_2$ and $c$ is a constant.

        \item Compute the multivariate MLE $[\hat{\lambda}, \hat{\theta}]$,
        and the observed information matrix $I(\hat{\lambda}, \hat{\theta})$.
    \end{enumerate}

    \begin{solution}
        \begin{enumerate}[label = (\alph*)]
            \item Represent $D_i \sim \Po(\lambda_i, Y_i)$,
            $i = 1, 2$,
            $D_i$ increases as $Y_i$ increases,
            or $\lambda_i$ increases.

            \begin{align*}
                \ell(\lambda_i) &= \frac{(y_i\lambda_i) ^ {d_i}e ^ {-\lambda_i y_i}}{d_i!}\footnotemark \\
                &= c_i\lambda_i ^ {d_i}e ^ {-\lambda_i y_i} \\
                \mathcal{L}(\lambda) &= \text{const} + d_i\log{\lambda_i} - \lambda_iy_i \\
                \mathcal{L};(\lambda) &= \frac{d_i}{\lambda_i} - y_i
                \intertext{At MLE,
                $\lambda_i = \hat{\lambda}_i$}
                0 &= \frac{d_i}{\hat{\lambda}_i} - y_i \iff \hat{\lambda}_i = \frac{d_i}{y_i}
            \end{align*}
            \footnotetext{Each sample is a Poisson data point.}
            \[
            \hat{\lambda}_1 = \frac{17}{2768.9} = 0.00614,\qquad \hat{\lambda}_2 = \frac{28}{1857.5} = 0.0151.
            \]

            Observed information:
            \[
            I(\hat{\lambda}_1) = -\mathcal{L}''(\hat{\lambda}_i) = -\left(-\frac{d_i}{\hat{\lambda}_i ^ 2}\right) = \frac{d_i}{(d_i / y_i) ^ 2} = \frac{y_i ^ 2}{d_i}.
            \]

            \item
            $\lambda = \lambda_1$ baseline rate $\theta = \frac{\lambda_2}{\lambda_1}$ multiplier to adjust for group $2$ which implies $\lambda_2 = \theta\lambda$.
            $\lambda = 1$ implies $\lambda_1 = \lambda_2$ the same,
            $\theta > 1$ implies $\lambda_2 > \lambda_1$ higher rate in group $2$.

            \[
            D_1 \sim \Po(\lambda y_i)\qquad D_2 \sim \Po(\theta\lambda y_2).
            \]

            \begin{align*}
                \ell(\lambda_1, \lambda_2) &= \frac{(\lambda_1 y_1) ^ {d_1}e ^ {-\lambda_1 y_1}}{d_1!} \times \frac{(\lambda_2y_2) ^ {d_2}e ^ {-\lambda_2y_2}}{d_2!} \\
                \ell(\lambda, \theta) &= \frac{(\lambda y_1) ^ {d_1}e ^ {-\lambda y_1}}{d_1!} \frac{(\theta\lambda y_2) ^ {d_2}e ^ {-\theta\lambda y_2)}}{d_2!} \\
                \mathcal{L}(\lambda, \theta) &= d_1\log(\lambda y_1) - \lambda y_1 - \log(d_1!) + d_2\log(\theta\lambda y_2) - \theta\lambda y_2 - \log(d_2!) \\
                &= \text{const} + d_1\log(\lambda) - \lambda y_1 + d_2\log(\theta) + d_2\log(\lambda) - \theta\lambda y_2 \\
                &= \text{const} + (d_1 + d_2)\log(\lambda) - \lambda y_1 - \theta\lambda y_2 + d_2\log(\theta).
            \end{align*}

            \item
            Find $\begin{bmatrix}
                \hat{\lambda} \\ \hat{\theta}
            \end{bmatrix}$.
            For $\hat{\lambda}$:
            \[
            \pd{\mathcal{L}}{\lambda} = \frac{d}{\lambda} - y_1 - \theta y_2\qquad\pd[\mathcal{L}]{\theta} = -\lambda y_2 + \frac{d_2}{\theta}.
            \]
            At MLE:
            \[
            \frac{d_2}{\hat{\theta}} = \hat{\lambda} y_2 \implies \hat{\lambda} = \frac{d_2}{\hat{\theta}y_2}
            \]
            \begin{align*}
                0 &= \frac{d}{\hat{\lambda}} - y_1 - \hat{\theta}y_2 \\
                &= \frac{d}{(d_2 / (\hat{\theta}y_2))} - y_1 - \hat{\theta}y_2 \\
                &= \frac{d\hat{\theta}y_2}{d_2} - y_1 - \hat{\theta}y_2 \\
                y_1 &= \hat{\theta}(((dy_2) / d_2) - y_2) \\
                \hat{\theta} &= \frac{y_1}{\frac{dy_2}{d_2} - y_2} = \dotsi = \frac{d_2 / y_2}{d_1 / y_1} = \frac{\hat{\lambda}_2}{\hat{\lambda}_1}
            \end{align*}
            ratio of group rates.
            \[
            \hat{\lambda} = \frac{d_2}{\hat{\theta}y_2} = \frac{d_2}{y_2}\frac{y_2d_1}{y_1d_2} = \frac{d_1}{d_1} = \hat{\lambda}_1.
            \]
        \end{enumerate}
    \end{solution}
\end{problem}

\newpage

\section{Problems Class \texorpdfstring{$2$}{}}

\begin{problem}[Sheet $5$ Question $8$]
    Suppose that $X_1, \dotsc, X_n$ form a random sample from a uniform distribution on the interval $[0, \theta]$,
    and that the following hypotheses are to be tested:
    \[
    H_0 : \theta \geq 2,\quad H_1 : \theta < 2.
    \]
    Let $Y_n = \max(X_1, \dotsc, X_n)$,
    and consider a test procedure such that the critical region contains all the outcomes for which $Y_n \leq 1.5$.

    \begin{enumerate}[label = (\alph*)]
        \item Determine the power function of the test.

        \item Determine the size of the test.
    \end{enumerate}

    \begin{solution}
        $X_1, \dotsc, X_n$ IID $U[0, \theta]$,
        $f(x; \theta) = \begin{cases}
            \frac{1}{\theta} & x \in [0, \theta] \\
            0 &\text{else}
        \end{cases}$.
        
        $H_0 : \theta \geq 2$,
        $H_1 : \theta < 2$.

        Test statistic:
        $Y_n = \max(X_1, \dotsc, X_n)$ and test based on $\mathcal{R} = \{\vec{x}\mid Y_n(x) \leq 1.5\}$,
        $I = \{y_n \mid Y_n(x) \leq 1.5\}$.
        
        \begin{enumerate}[label = (\alph*)]
            \item
            Power function is $\pi(\theta) = \P(Y_n(x) \leq 1.5; \theta)$.
            Need the distribution for $Y_n$
            (CDF)
            \begin{align*}
                \P(Y_n \leq y) &= \P(\max(X_1, \dotsc, X_n) \leq y) \\
                &= \P(X_1 \leq y, X_2 \leq y, \dotsc, X_n \leq y) \\
                &= \P(X_1 \leq y)\P(X_2 \leq y)\dotsi \P(X_n \leq y) \\
                &= \prod_{i = 1}^{n}\P(X_i \leq y) \\
                &= \P(X \leq y) ^ n.
            \end{align*}
            What is $\P(X \leq y)$?
            For $y \in [0, \theta]$,
            \[
            \P(X \leq y) = \int_{0}^{y}\frac{1}{\theta}\,dx = \left[\frac{x}{\theta}\right]_{0}^{y} = \frac{y}{\theta}
            \]
            \[
            \P(Y_n \leq y) = \begin{cases}
                0 & y < 0 \\
                \left(\frac{y}{\theta}\right) ^ n & y \in [0, \theta] \\
                1 & y > 0.
            \end{cases}
            \]
            Test rejects for any $Y_n \leq 1.5$,
            so if $\theta \leq 1.5$,
            then all $X_i \leq 1.5$ which implies $Y_n \leq 1.5$ and we always reject.

            Conversely,
            for $\theta > 1.5$,
            the probability of rejection is $\P(Y_n \leq 1.5) = \left(\frac{1.5}{\theta}\right) ^ n$ so
            \[
            \pi(\theta) = \begin{cases}
                1 & \theta \leq 1.5 \\
                \left(\frac{1.5}{\theta}\right) ^ n & \theta > 1.5.
            \end{cases}
            \]

            \item
            Size of the test.
            The significance level of the test is $\pi(\theta) = \P(\text{Type I})$ under $H_0$,
            which changes with $\theta$.
            The size is max/sup of the level/significance over all $\theta \in \Omega_0$
            \[
            \alpha = \sup_{\theta \geq 2}\pi(\theta) = \sup_{\theta \geq 2}\left(\frac{1.5}{\theta}\right) ^ n = \left(\frac{3}{4}\right) ^ n
            \]
        \end{enumerate}
    \end{solution}
\end{problem}

\begin{problem}[Sheet $4$ Question $4$]
    A simple model for the drug concentration in plasma over time after a single intravenous injection is $c(t) = \theta_2\exp\{-\theta_1t\}$,
    with $\theta_1, \theta_2 > 0$.
    For simplicity,
    we assume here that $\theta_2 = 1$.
    
    \begin{enumerate}[label = (\alph*)]
        \item Assume that $n$ plasma samples had their concentrations $c_i$,
        $i = 1, \dotsc, n$,
        measured at the same single time-point $t$ and assume that the $c_i$ are IID $N(c(t), \sigma ^ 2)$ is appropriate for the data.
        Calculate the MLE of $\theta_1$.

        \item Calculate the asymptotic variance of the MLE.

        \item In pharmacokinetic studies one is often interested in the area under the concentration curve
        \[
        \alpha = \int_{0}^{\infty}\exp\{-\theta_1t\}\,dt.
        \]
        Calculate the MLE for $\alpha$,
        and its variance estimate in terms of $\hat{\theta}_1$ by using the delta theorem.

       \item We now would like to determine the optimal time point for measuring the concentrations $c_i$.
       It was suggested to identify the time point for which the asymptotic variance of the MLE of $1$ is smallest.
       Assume that $1$ is known and minimise this variance with respect to $t$ to find the optimal time point $t_{opt}$.
    \end{enumerate}

    \begin{solution}\phantom{}
        \begin{enumerate}[label = (\alph*)]
            \item
            Normal model for $c_i$,
            with mean
            \[
            c_i \sim N(e ^ {-\theta_1t}, \sigma ^ 2)
            \]
            $\theta_1 > 0$,
            assume $\sigma ^ 2$ known.
            \[
            f(c_i; t, \sigma ^ 2, \theta_1) = \frac{1}{\sqrt{2\pi\sigma ^ 2}}\exp\left\{-\frac{1}{2\sigma ^ 2}(c_i - e ^ {-\theta_1t}) ^ 2\right\}
            \]
            implies
            \begin{align*}
                \ell(\theta_1) &= (2\pi\sigma ^ 2) ^ {-n / 2}\exp\left\{-\frac{1}{2\sigma ^ 2}(c_i - e ^ {-\theta_1t}) ^ 2\right\} \\
                \mathcal{L}(\theta_1) &= \text{const} - \frac{1}{2\sigma ^ 2}\sum_{i = 1}^{n}(c_i - e ^ {-\theta_1t}) ^ 2 \\
                \mathcal{L}'(\theta_1) &= -\frac{1}{2\sigma ^ 2}\sum_{i = 1}^{n}2(c_i - e ^ {-\theta_1t})\times te ^ {-\theta_1 t} \\
                &= -\frac{te ^ {-\theta_1t}}{\sigma ^ 2}\sum_{i}(c_i - e ^ {-\theta_1t}) \\
                &= -\frac{te ^ {-\theta_1t}}{\sigma ^ 2}(n\bar{c} - ne ^ {-\theta_1t}) \\
                &= \frac{nte ^ {-\theta_1t}}{\sigma ^ 2}(e ^ {-\theta_1t} - \bar{c})
                \intertext{Let $g(\theta_1) = e ^ {-\theta_1t}$}
                &= \frac{ntg(\theta_1)}{\sigma ^ 2}(g(\theta_1) - \bar{c}).
            \end{align*}
            At the MLE:
            \[
            0 = \frac{ntg(\hat{\theta}_1)}{\sigma ^ 2}(g(\hat{\theta}_1) - \bar{c})
            \]
            only zero when $g(\hat{\theta}_1) = \bar{c}$,
            \[
            e ^ {-\hat{\theta}_1t} = \bar{c} \iff \hat{\theta}_1 = -\frac{\log{\bar{c}}}{t}.
            \]
            To verify maximum,
            differentiate again:
            \begin{align*}
                \mathcal{L}''(\theta) &= \pd{\theta_1}\frac{nt}{\sigma ^ 2}g(\theta_1)(g(\theta_1) - \bar{c})) \\
                &= \frac{nt}{\sigma ^ 2}(2g(\theta_1)g'(\theta_1) - \bar{c}g'(\theta_1)) \\
                &= -\frac{nt ^ 2g(\theta_1)}{\sigma ^ 2}(2g(\theta_1) - \bar{c}).
            \end{align*}
            Hard to show less than $0$ for all $\theta$,
            look at $\hat{\theta}$.
            \[
            \mathcal{L}''(\hat{\theta}) = -\frac{nt ^ 2\bar{c} ^ 2}{\sigma ^ 2} < 0.
            \]

            \item
            Asymptotic
            (large-sample,
            $n$ large)
            variance is the Crámer-Rao lower bound:
            \[
            \frac{1}{\mathbb{I}_n(\theta_1)}.
            \]
            We have $\mathcal{L}''(\theta)$,
            so
            \begin{align*}
                \mathbb{I}_n(\theta) &= \E_C(-\mathcal{L}''(\theta)) \\
                &= \frac{nt}{\sigma ^ 2}g(\theta_1)(2g(\theta_1) - \E(\bar{C}))
                \intertext{As $\E(\bar{C}) = \E(C) = e ^ {-\theta_1t} = g(\theta_1)$.}
                \mathbb{I}_n(\theta) &= \frac{nt ^ 2}{\sigma ^ 2}g(\theta_1)(2g(\theta_1) - g(\theta_1)) \\
                &= \frac{nt ^ 2}{\sigma ^ 2}g(\theta_1) ^ 2 \\
                &= \frac{nt ^ 2}{\sigma ^ 2}\exp\{-2\theta t\}.
            \end{align*}
            For large $n$:
            $\Var(\hat{\theta}) = \frac{\sigma ^ 2}{n}\frac{1}{t ^ 2}\exp\{2\theta_1t\}$.

            \item Want 
            \begin{align*}
                \alpha &= \int_{0}^{\infty}\exp\{-\theta_1t\}\,dt \\
                &= \left[\frac{-\exp\{-\theta_1t\}}{\theta_1}\right]_{0}^{\infty} \\
                &= \frac{1}{\theta_1} = \alpha(\theta_1).
            \end{align*}
            MLE of $\alpha$:
            \[
            \hat{\alpha} = \frac{1}{\hat{\theta}_1} = -\frac{t}{\log{\bar{c}}}.
            \]
            \[
            \Var(\hat{\alpha}) = \Var(\hat{\theta}_1)\left(\pd[\alpha]{\theta_1}\right) ^ 2 = \left(\frac{\sigma ^ 2}{nt ^ 2}\exp\{2\theta_1 t\}\right)\left(-\frac{1}{\theta_1}\right) ^ 2 = \frac{\sigma ^ 2}{nt ^ 2\theta_1 ^ 4}\exp\{2\theta_1t\}.
            \]
            CI for $\alpha$:
            \[
            -\frac{t}{\log{\bar{c}}} \pm z_{\alpha / 2}^{*}\frac{\sigma}{t\hat{\theta_1} ^ 2\sqrt{n}}\exp\{\theta_1 t\}.
            \]
        \end{enumerate}
    \end{solution}
\end{problem}



















\end{document}